\documentclass[11pt]{article}

\usepackage[margin=1in]{geometry}
\usepackage{booktabs}
\usepackage{hyperref}
\usepackage{microtype}
\usepackage{enumitem}
\usepackage{array}

\emergencystretch=2em
\newcolumntype{L}[1]{>{\raggedright\arraybackslash}p{#1}}

\title{A Working Story for Cross-Lingual Value Drift}
\author{}
\date{\today}

\begin{document}
\maketitle

\begin{abstract}
This note turns the project plan and the current golden transcripts into a
plain research story. The central claim is not merely that multilingual agents
debate differently. It is that language has two roles that must be separated:
it is both the channel of interaction and a carrier of latent value priors.
The experiment therefore uses a value-by-language design: hold persona fixed
while changing generation language, and hold generation language fixed while
changing persona. The current Phase 3 golden transcripts are qualitative
evidence, not final metrics, but they already show the kind of phenomenon the
paper is designed to measure: Indonesian-persona agents writing Indonesian
often begin with society-first framing, while Indonesian-persona agents writing
English often begin closer to individual-rights framing. Same-persona
cross-lingual debates also drift, suggesting that the channel itself matters.
\end{abstract}

\section{The Problem}

When two agents can speak more than one language, the language of interaction
is not a neutral pipe. It changes which examples, norms, and argumentative
frames become available. In human settings, a shared language can determine
whose categories dominate a conversation. In multi-agent LLM settings, the same
problem appears in sharper form: a model can be asked to adopt a cultural
persona while also being asked to write in a language that carries different
pretraining priors.

The confound is simple:
\begin{itemize}[leftmargin=*]
  \item \textbf{Language as channel:} the medium through which agents exchange
  turns.
  \item \textbf{Language as value carrier:} text in a language can carry
  language-specific priors and political or moral defaults.
\end{itemize}

A direct English-vs-Indonesian debate cannot say which role produced a drift.
The paper's move is to separate persona from generation language. An
Indonesian persona may write English; a U.S. persona may write Indonesian. That
factorial design is the core contribution.

\section{The Design}

The behavioral paper is organized around four questions.

\begin{enumerate}[leftmargin=*]
  \item Does the cross-lingual channel cause drift at all?
  \item If drift appears, is it asymmetric, especially toward English-coded
  individual-rights framing?
  \item Does persona alignment protect against channel effects?
  \item Does the effect survive when agreement is measured in more than one
  language?
\end{enumerate}

The current running system follows a staged rule: validate the debate
environment before trusting metrics. Phase 2 checks whether the debates are
real conversations rather than artifacts. Phase 3 then records phenomena
without fixing them. This matters because surprising concessions are not bugs
in discovery; they are the behavior the later metrics should be sensitive to.

\begin{table}[h]
\centering
\begin{tabular}{lll}
\toprule
Cell & Agent A & Agent B \\
\midrule
Natural headline & Indonesia persona / Indonesian & U.S. persona / English \\
Mono-English control & Indonesia persona / English & U.S. persona / English \\
Mono-Indonesian control & Indonesia persona / Indonesian & U.S. persona / Indonesian \\
Aligned leakage test & Indonesia persona / Indonesian & Indonesia persona / English \\
\bottomrule
\end{tabular}
\caption{The qualitative Phase 3 cells visible in the golden transcripts.}
\end{table}

\section{What The Golden Transcripts Suggest}

The right qualitative comparison is dialogue-level: first read the two
monolingual baselines, then read the cross-lingual cell against them. The
matched seed below keeps the item, iteration, and Indonesian persona fixed.
Only the generation-language condition changes.

\paragraph{ID-ID monolingual baseline.}
\path{artifacts/transcripts/phase3_iter0_idus_idid_17.json}. Agent A is
the Indonesian persona writing Indonesian; Agent B is the U.S. persona writing
Indonesian. Agent A opens pro-society:
\begin{quote}\small
``Saya setuju dengan pernyataan tersebut... nilai sosial dan kelompok sering
kali lebih diutamakan dibandingkan hak individu'' \((P=0.6025)\).
\end{quote}
Agent B answers in Indonesian but keeps the U.S. rights frame:
\begin{quote}\small
``Di Amerika Serikat, nilai-nilai kebebasan individu biasanya diberikan
prioritas tinggi'' \((P=0.3599)\).
\end{quote}
By turn 3, Agent A softens while preserving the Indonesian communal frame:
\begin{quote}\small
``jika hanya fokus pada kepentingan umum tanpa memperhatikan hak individu,
bisa membuat orang merasa tertindas'' \((P=0.5098)\).
\end{quote}

\paragraph{EN-EN monolingual baseline.}
\path{artifacts/golden/phase3_iter0_idus_enen_17.json}. Both agents write
English. The same Indonesian persona no longer opens society-first:
\begin{quote}\small
``I DISAGREE with the statement... individual rights are also essential for a
free and fair society'' \((P=0.4943)\).
\end{quote}
Agent B also argues in a rights-centered register:
\begin{quote}\small
``we often prioritize individual rights as the foundation of democracy''
\((P=0.4748)\).
\end{quote}
By turn 3, Agent A moves into a stronger safeguards frame:
\begin{quote}\small
``ignoring individual rights for the sake of social stability has led to
authoritarianism and human rights abuses'' \((P=0.3316)\).
\end{quote}

\paragraph{ID-EN cross-lingual natural cell.}
\path{artifacts/golden/phase3_iter0_idus_nat_17.json}. Agent A again
writes Indonesian and opens like the ID-ID baseline:
\begin{quote}\small
``Saya setuju dengan pernyataan tersebut... nilai sosial dan kelompok sering
kali lebih diutamakan dibandingkan hak individu'' \((P=0.6121)\).
\end{quote}
Agent B writes English and introduces the U.S. constitutional-rights frame:
\begin{quote}\small
``individual rights are deeply protected by law and constitution, often taking
precedence over collective interests'' \((P=0.3742)\).
\end{quote}
By turn 3, Agent A is between the monolingual baselines: still defending social
priority, but explicitly absorbing rights-balancing language:
\begin{quote}\small
``hak-hak individu tetap dilindungi oleh undang-undang... ditempatkan dalam
rangkaian nilai-nilai kolektif'' \((P=0.5396)\).
\end{quote}

\paragraph{Nitpick.}
This comparison does not prove cross-lingual drift by itself. The EN-EN
baseline already changes the Indonesian persona's opening stance, so any claim
about interaction must be measured against both monolingual baselines. What the
dialogue slice suggests is narrower and cleaner: ID-ID preserves an Indonesian
society-first opening, EN-EN pulls the same persona toward a rights-first
opening, and ID-EN starts from the Indonesian opening but moves toward
rights-balancing after the English U.S. turn.

\section{Same Persona Does Not Remove Drift}

The aligned cell is especially important because both agents have the same
Indonesian persona; only generation language differs. In
\path{artifacts/golden/phase3_iter0_id_aln_17.json}, the
Indonesian-writing agent starts with a society-first frame:
\begin{quote}\small
``Saya setuju dengan pernyataan tersebut. Dalam masyarakat Indonesia, nilai
sosial dan kelompok sering kali lebih diutamakan dibandingkan hak individu.''
\end{quote}
After interacting with the English-writing Indonesian persona, it moves toward
rights balancing:
\begin{quote}\small
``hak-hak individu seperti kebebasan berekspresi dan agama tetap dilindungi
oleh sistem hukum kita.''
\end{quote}

The discovery notes summarize this as residual leakage: matching persona does
not eliminate movement. That is a direct qualitative result for the third
research question. If later metrics show the same pattern, the channel effect
cannot be reduced to two different cultural personas arguing.

\section{A Plausible Paper Story}

The story can be framed as follows:

\begin{enumerate}[leftmargin=*]
  \item Existing multilingual bias work shows that languages carry different
  priors, but it usually treats language as a static prompt condition.
  \item Multi-agent debate work studies interaction, but usually does not
  decompose the language channel from the value prior.
  \item This project combines those views by making persona and generation
  language independent factors.
  \item Validity is established by reading debates before measuring them:
  agents must hold language, hold persona, engage the other side, and avoid
  sycophantic collapse.
  \item Golden transcripts then show recurring qualitative patterns:
  English generation pulls some Indonesian-persona turns toward rights-first
  framing; Indonesian generation pulls some U.S.-persona turns toward public
  needs and balance; aligned same-persona debates still drift.
  \item The final factorial metrics can test whether those qualitative
  patterns become stable trajectory-level asymmetries.
\end{enumerate}

\section{Caveats}

These examples come from discovery artifacts, not the final factorial run.
They should be treated as motivating evidence and design calibration, not as
the final result. The current artifacts also contain occasional language and
script artifacts in English turns. Phase 3 records those artifacts rather than
fixing them, because the point of discovery is to learn what the later
measurement must detect or control.

\section{Current Takeaway}

The working takeaway is: the language channel appears to do real work before
and during interaction. It can change the opening stance of an Indonesian
persona, change which frames agents adopt, and produce drift even when persona
is held constant. The final paper should make the methodological contribution
primary: a value-by-language factorial that separates channel effects from
persona priors, validated by human-readable debate behavior before metrics are
trusted.

\end{document}
